% ---------------------------------------------------------
% Definition --- Differential Equation
% ---------------------------------------------------------

\begin{definitionbox}{Definition}
\begin{definition}[Differential Equation]
\label{def:ordinary-differential-equations-terminology-differential-equation}
A \textbf{differential equation} is an equation whose unknown is a function and whose terms involve one or more derivatives of that function.

For an unknown function $y=y(t)$, a differential equation may be written in the form
\[
F(t,y,y',y'',\dots,y^{(n)})=0.
\]
Here $t$ is the independent variable, $y(t)$ is the unknown function, and
\[
y',y'',\dots,y^{(n)}
\]
are derivatives of $y$.
\end{definition}
\end{definitionbox}

\begin{remark*}[Standard quantified statement]
This definition fixes the named kind of differential-equation object by the
conditions stated in the definition.
\end{remark*}

\begin{remark*}[Interpretation]
At the computational level, a differential equation describes how a quantity changes. Instead of giving the function $y(t)$ directly, it gives a relation involving the rate of change of $y(t)$, and possibly higher-order rates of change.
\end{remark*}

\begin{remark*}[Examples]
The equation
\[
y'=3y
\]
is a differential equation because it relates the unknown function $y$ to its derivative $y'$.

A solution is
\[
y(t)=Ce^{3t},
\]
where $C$ is a constant. Indeed,
\[
y'(t)=3Ce^{3t}=3y(t).
\]
\end{remark*}

\begin{dependencies}
\begin{itemize}
  \item \hyperref[def:function]{Function}
  \item \hyperref[def:function-equality]{Function Equality}
  \item \hyperref[def:derivative-at-a-point]{Derivative at a Point}
\end{itemize}
\end{dependencies}

% ---------------------------------------------------------
% Definition --- Order of a Differential Equation
% ---------------------------------------------------------

\begin{definitionbox}{Definition}
\begin{definition}[Order of a Differential Equation]
\label{def:ordinary-differential-equations-terminology-order-of-a-differential-equation}
The \textbf{order} of a differential equation is the highest derivative of the unknown function that appears in the equation.

For example, if an equation involving an unknown function $y=y(t)$ can be written in the form
\[
F(t,y,y',y'',\dots,y^{(n)})=0
\]
and $y^{(n)}$ is the highest derivative that appears, then the differential equation has \textbf{order} $n$.
\end{definition}
\end{definitionbox}

\begin{remark*}[Standard quantified statement]
This definition fixes the named kind of differential-equation object by the
conditions stated in the definition.
\end{remark*}

\begin{remark*}[Interpretation]
The order records the highest derivative that controls the type of differential
equation being studied.
\end{remark*}

\begin{remark*}[Examples]
The equation
\[
y'=3y
\]
is first order, since the highest derivative appearing is $y'$.

The equation
\[
y''+4y'+5y=0
\]
is second order, since the highest derivative appearing is $y''$.

The equation
\[
y^{(4)}-t y''+y=0
\]
is fourth order, since the highest derivative appearing is $y^{(4)}$.
\end{remark*}

\begin{dependencies}
\begin{itemize}
  \item \hyperref[def:ordinary-differential-equations-terminology-differential-equation]{Differential Equation}
  \item \hyperref[def:derivative-at-a-point]{Derivative at a Point}
\end{itemize}
\end{dependencies}

% ---------------------------------------------------------
% Definition --- First-Order Differential Equation
% ---------------------------------------------------------

\begin{definitionbox}{Definition}
\begin{definition}[First-Order Differential Equation]
\label{def:ordinary-differential-equations-terminology-first-order-differential-equation}
A \textbf{first-order differential equation} is a differential equation whose
order is one.

Equivalently, it is a differential equation involving the first derivative of
the unknown function but no derivative of higher order.
\end{definition}
\end{definitionbox}

\begin{remark*}[Standard quantified statement]
This definition fixes the named differential-equation class by the order
condition stated in the definition.
\end{remark*}

\begin{remark*}[Interpretation]
First-order equations are the basic setting for many evolution models because
they relate the current state of a quantity to its instantaneous rate of
change.
\end{remark*}

\begin{dependencies}
\begin{itemize}
  \item \hyperref[def:ordinary-differential-equations-terminology-differential-equation]{Differential Equation}
  \item \hyperref[def:ordinary-differential-equations-terminology-order-of-a-differential-equation]{Order of a Differential Equation}
\end{itemize}
\end{dependencies}

% ---------------------------------------------------------
% Definition --- Solution of a Differential Equation
% ---------------------------------------------------------

\begin{definitionbox}{Definition}
\begin{definition}[Solution of a Differential Equation]
\label{def:ordinary-differential-equations-terminology-solution-of-a-differential-equation}
A \textbf{solution} of a differential equation is a function that satisfies the equation on a specified interval.

More precisely, if a differential equation is written in the form
\[
F(t,y,y',y'',\dots,y^{(n)})=0,
\]
then a solution on an interval $I$ is a function $y:I\to\mathbb{R}$ such that $y$ has all derivatives required by the equation and
\[
F(t,y(t),y'(t),y''(t),\dots,y^{(n)}(t))=0
\]
for every $t\in I$.
\end{definition}
\end{definitionbox}

\begin{remark*}[Standard quantified statement]
This definition fixes the named kind of differential-equation object by the
conditions stated in the definition.
\end{remark*}

\begin{remark*}[Interpretation]
A solution is not a number; it is a function. The function solves the differential equation when substituting the function and its derivatives into the equation makes the equation true for every point in the interval under consideration.
\end{remark*}

\begin{remark*}[Examples]
The function
\[
y(t)=Ce^{3t}
\]
is a solution of
\[
y'=3y
\]
on $\mathbb{R}$, because
\[
y'(t)=3Ce^{3t}=3y(t)
\]
for every $t\in\mathbb{R}$.
\end{remark*}

\begin{dependencies}
\begin{itemize}
  \item \hyperref[def:ordinary-differential-equations-terminology-differential-equation]{Differential Equation}
  \item \hyperref[def:function]{Function}
  \item \hyperref[def:interval]{Interval}
\end{itemize}
\end{dependencies}

% ---------------------------------------------------------
% Definition --- Autonomous Differential Equation
% ---------------------------------------------------------

\begin{definitionbox}{Definition}
\begin{definition}[Autonomous Differential Equation]
\label{def:ordinary-differential-equations-terminology-autonomous-differential-equation}
An \textbf{autonomous differential equation} is a differential equation in which the independent variable does not appear explicitly.

For a first-order ordinary differential equation, an autonomous equation has the form
\[
y'=f(y),
\]
where the rate of change of $y$ depends only on the current value of $y$, not directly on time $t$.
\end{definition}
\end{definitionbox}

\begin{remark*}[Standard quantified statement]
This definition fixes the named kind of differential-equation object by the
conditions stated in the definition.
\end{remark*}

\begin{remark*}[Interpretation]
An autonomous differential equation describes a system whose rule of evolution does not change with time. If the system is in the same state $y$ at two different times, then the differential equation assigns the same rate of change at both times.
\end{remark*}

\begin{remark*}[Examples]
The equation
\[
y'=3y
\]
is autonomous, since the right-hand side depends only on $y$.

The equation
\[
y'=t+y
\]
is not autonomous, since the independent variable $t$ appears explicitly.
\end{remark*}

\begin{dependencies}
\begin{itemize}
  \item \hyperref[def:ordinary-differential-equations-terminology-first-order-differential-equation]{First-Order Differential Equation}
  \item \hyperref[def:function]{Function}
\end{itemize}
\end{dependencies}

% ---------------------------------------------------------
% Definition --- Non-Autonomous Differential Equation
% ---------------------------------------------------------

\begin{definitionbox}{Definition}
\begin{definition}[Non-Autonomous Differential Equation]
\label{def:ordinary-differential-equations-terminology-non-autonomous-differential-equation}
A \textbf{non-autonomous differential equation} is a differential equation in which the independent variable appears explicitly.

For a first-order ordinary differential equation, a non-autonomous equation has the form
\[
y'=f(t,y),
\]
where the rate of change of $y$ may depend both on the current value of $y$ and on the independent variable $t$.
\end{definition}
\end{definitionbox}

\begin{remark*}[Standard quantified statement]
This definition fixes the named kind of differential-equation object by the
conditions stated in the definition.
\end{remark*}

\begin{remark*}[Interpretation]
A non-autonomous differential equation describes a system whose rule of evolution may change over time. Even if the system has the same state $y$ at two different times, the differential equation may assign different rates of change because the value of $t$ may be different.
\end{remark*}

\begin{remark*}[Examples]
The equation
\[
y'=t+y
\]
is non-autonomous, since the independent variable $t$ appears explicitly.

The equation
\[
y'=3y
\]
is autonomous, not non-autonomous, since the right-hand side depends only on $y$.
\end{remark*}

\begin{dependencies}
\begin{itemize}
  \item \hyperref[def:ordinary-differential-equations-terminology-first-order-differential-equation]{First-Order Differential Equation}
  \item \hyperref[def:function]{Function}
\end{itemize}
\end{dependencies}

% ---------------------------------------------------------
% Definition --- Linear Differential Equation
% ---------------------------------------------------------

\begin{definitionbox}{Definition}
\begin{definition}[Linear Differential Equation]
\label{def:ordinary-differential-equations-terminology-linear-differential-equation}
A differential equation is called \textbf{linear} if the unknown function and its derivatives appear only to the first power and are not multiplied together.

An $n$th-order linear ordinary differential equation for an unknown function $y=y(t)$ has the form
\[
a_n(t)y^{(n)}+a_{n-1}(t)y^{(n-1)}+\cdots+a_1(t)y'+a_0(t)y=g(t),
\]
where the coefficient functions
\[
a_0(t),a_1(t),\dots,a_n(t)
\]
and the forcing term $g(t)$ are given functions of the independent variable $t$.
\end{definition}
\end{definitionbox}

\begin{remark*}[Standard quantified statement]
This definition fixes the named kind of differential-equation object by the
conditions stated in the definition.
\end{remark*}

\begin{remark*}[Interpretation]
Linearity means that the equation is linear in the unknown function $y$ and its derivatives. The coefficients may depend on $t$, but they may not depend on $y$ or its derivatives. Thus terms such as $y^2$, $yy'$, $\sin(y)$, and $(y')^2$ make an equation nonlinear.
\end{remark*}

\begin{remark*}[Examples]
The equation
\[
y''+3t y'+2y=\sin t
\]
is linear.

The equation
\[
y''+y^2=0
\]
is nonlinear because it contains the term $y^2$.

The equation
\[
y'+yy'=t
\]
is nonlinear because it contains the product $yy'$.
\end{remark*}

\begin{dependencies}
\begin{itemize}
  \item \hyperref[def:ordinary-differential-equations-terminology-differential-equation]{Differential Equation}
  \item \hyperref[def:linear-combination-of-real-valued-functions]{Linear Combination of Real-Valued Functions}
  \item \hyperref[def:pointwise-scalar-multiple-of-a-function]{Pointwise Scalar Multiple of a Function}
\end{itemize}
\end{dependencies}

% ---------------------------------------------------------
% Definition --- Homogeneous Linear Differential Equation
% ---------------------------------------------------------

\begin{definitionbox}{Definition}
\begin{definition}[Homogeneous Linear Differential Equation]
\label{def:ordinary-differential-equations-terminology-homogeneous-linear-differential-equation}
A \textbf{homogeneous linear differential equation} is a linear differential
equation whose forcing term is zero.

For an $n$th-order linear ordinary differential equation
\[
a_n(t)y^{(n)}+a_{n-1}(t)y^{(n-1)}+\cdots+a_1(t)y'+a_0(t)y=g(t),
\]
homogeneity means that
\[
g(t)=0
\]
on the interval under consideration.
\end{definition}
\end{definitionbox}

\begin{remark*}[Standard quantified statement]
This definition fixes the named differential-equation class by the zero forcing
condition stated in the definition.
\end{remark*}

\begin{remark*}[Interpretation]
The homogeneous case records the intrinsic dynamics of the linear equation,
without an external forcing term. Nonzero forcing terms produce nonhomogeneous
linear differential equations.
\end{remark*}

\begin{dependencies}
\begin{itemize}
  \item \hyperref[def:ordinary-differential-equations-terminology-linear-differential-equation]{Linear Differential Equation}
  \item \hyperref[def:function-equality]{Function Equality}
\end{itemize}
\end{dependencies}
